\section{結言}
今回のDaily Xでは、モデル能力の競争と並行して、生成AIを実運用へ組み込むための周辺層が前面に出た。OpenAIのPrivate Safety Processingはデータ保持と安全監視、Grok 4.6のBedrock展開はクラウド流通、AgentCore Paymentsはエージェントによる取引という、それぞれ異なる実装上の境界を扱っている。

一方、Claudeによるタンパク質設計は、汎用モデルが科学ワークフローへ深く入る方向を示した。QwenやDeepSeekの話題ではローカル実行・コスト・実利用性能への関心が続き、公式発表よりコミュニティ検証が先行する場面も多い。X上では、生成AIの評価軸が単純なモデル性能だけでなく、安全、配備、決済、科学応用へ拡張している一日だった。
